\documentclass[11pt,a4paper]{exam}
\noprintanswers

\usepackage[T1]{fontenc}
\usepackage{lmodern}
\usepackage{amsmath,amssymb}
\usepackage[margin=2cm,includeheadfoot,headheight=15pt,headsep=18pt]{geometry}
\usepackage{needspace}
\usepackage[hidelinks]{hyperref}

\newcommand{\CourseTitle}{Introduction to Koszul Duality}
\newcommand{\AssignmentNumber}{1}
\newcommand{\Deadline}{to be announced}

\pagestyle{headandfoot}
\header{\small\CourseTitle}
       {\small Assignment \AssignmentNumber}
       {\small Deadline: \Deadline}
\headrule
\footer{}{\small Submit via email}{\small\thepage}

\newcommand{\kk}{\Bbbk}
\newcommand{\bbD}{\mathbb D}
\newcommand{\grstar}{(*)}
\DeclareMathOperator{\Hom}{Hom}
\DeclareMathOperator{\End}{End}
\DeclareMathOperator{\rad}{rad}
\DeclareMathOperator{\Sym}{Sym}
\newcommand{\op}{\mathrm{op}}
\newcommand{\id}{\mathrm{id}}

\qformat{\textbf{Exercise \thequestion.\ \thequestiontitle}\hfill}
\renewcommand{\questionshook}{%
  \setlength{\leftmargin}{0pt}%
  \setlength{\labelwidth}{-\labelsep}%
  \setlength{\itemsep}{13pt}%
}
\renewcommand{\thepartno}{\arabic{partno}}
\renewcommand{\partlabel}{(\thepartno)}
\renewcommand{\partshook}{\setlength{\itemsep}{4pt}}
\setlength{\parindent}{0pt}
\setlength{\parskip}{4pt}
\setlength{\emergencystretch}{1em}
\raggedbottom

% Revision notes (v2), not printed:
% * Every part of every question is meant to be accessible; no part is
%   labelled "advanced".  Optional material is labelled "Bonus".
% * Suggested selection rule: at least two full questions, at least one of
%   which is Exercise 5 or 6 (Lecture 3).
% * Hints are labelled and kept apart from the questions.

\begin{document}

\textbf{Exercises from Lectures 1--3.}
Throughout, $\kk$ is a field, and all modules are right modules unless a
left action is specified. Write $\kk A_n$ for the algebra of upper triangular
$n\times n$ matrices over $\kk$, with matrix units $e_{ij}$ ($1\leq i\leq j\leq n$),
$e_{ij}e_{pq}=\delta_{jp}e_{iq}$, and grading $\deg e_{ij}=j-i$.
For a graded module $M$, the shift convention is $M(r)_d=M_{d+r}$.
Ideals in algebra presentations are two-sided ideals.
For $\kk A_n$ put $P_i=e_{ii}\kk A_n$ (graded by the inherited grading), and
write $\Hom^d_A(M,N)$ for the homogeneous $A$-linear maps $M\to N$ of degree $d$.

\textbf{Hilbert series.}
For a graded algebra $A=\bigoplus_{d\geq0}A_d$ with
$A_0=\bigoplus_{i=1}^r\kk e_i$ (coordinate idempotents, in the listed order),
\[
  h_A(t)_{ij}=\sum_{d\geq0}\dim_{\kk}(e_iA_de_j)\,t^d\qquad(1\leq i,j\leq r),
\]
so the coefficient of $t^d$ is the dimension vector of the $A_0$-bimodule
$A_d$, recorded with two indices. When $r=1$ there is a single entry.
Arrays are multiplied as matrices: $(CD)_{i\ell}=\sum_jc_{ij}d_{j\ell}$, with
the same rule for series.

\textbf{Quadratic duals.}
For $A=T_S(V)/(R)$ with $R\subset V\otimes_SV$, put $V^*=\Hom_S(V,S)$ (in
degree $1$), and let
\[
  \theta_r(g\otimes f)(v\otimes w)=g(f(v)w)\qquad(f,g\in V^*).
\]
The \emph{orthogonal complement} of $R$ is
$R^\perp=\{F\in V^*\otimes_SV^*\mid\theta_r(F)(r)=0\text{ for all }r\in R\}$,
and $A^!=T_S(V^*)/(R^\perp)$.
When $S=\bigoplus_i\kk e_i$ and $v\in e_iVe_j$ is a chosen basis vector,
its dual $v^*:V\to S$ takes $v$ to $e_j$ and the other chosen basis vectors
to zero; it satisfies $e_jv^*e_i=v^*$, and these $v^*$ form a basis of $V^*$.
In particular $\dim_\kk e_jV^*e_i=\dim_\kk e_iVe_j$.

\medskip
\hrule

\section*{Lecture 1: Graded algebras, modules, and duals}
\begin{questions}

\titledquestion{Projectives, endomorphisms, and duals over $\kk A_3$}
Let $A=\kk A_3$. The parts are independent except where stated.
\begin{parts}
\part Verify that $A_d=\operatorname{span}_\kk\{e_{ij}\mid j-i=d\}$ is a
grading of $\kk A_n$ for every $n\geq1$, and compute $\rad(\kk A_3)$ from its
definition as the intersection of all maximal right ideals.

\emph{Hint.} Let $J$ be the strictly upper triangular matrices; then $J^3=0$.
If $L$ is a simple right module, then $LJ\in\{0,L\}$, and $LJ=L$ would give
$L=LJ^3=0$. Then identify the maximal right ideals of $A/J\cong\kk^{\oplus3}$.

\part For each of the nine ordered pairs $(i,j)$, find a basis of
$\Hom_A(P_i,P_j)$ and the degree of each basis map, in a $3\times3$ table
(rows: source $P_i$; columns: target $P_j$).
Then determine, for all $i,j$ and $r,s\in\mathbb Z$, when there is a non-zero
degree-preserving map $P_i(r)\to P_j(s)$.

\emph{Hint.} An $A$-linear map $u:P_i\to P_j$ is determined by $u(e_{ii})$,
and $u(e_{ii})=u(e_{ii})e_{ii}$.

\part Let $\lambda:A\to\End_A(A_A)$ send $a$ to $x\mapsto ax$.
\begin{subparts}
\subpart Show that $\lambda$ is a bijective algebra homomorphism, and that
$\lambda(e_{ji})$ restricts to the basis map $P_i\to P_j$ of part (2).
\subpart Prove that
\[
  h_A(t)_{ij}=\sum_{d\in\mathbb Z}\dim_\kk\Hom^d_A(P_j,P_i)\,t^d
  \qquad(1\leq i,j\leq3),
\]
and compare with the matrix $h_{\kk A_3}(t)$ from Lecture~2.
\end{subparts}

\emph{Hint for (ii).} Evaluate a map at $e_{jj}$.

\part Let $M=e_{22}A$ and
$\bbD(M)_d=\Hom_\kk(M_{-d},\kk)$ with $(af)(m)=f(ma)$.
\begin{subparts}
\subpart Find a homogeneous basis of $\bbD(M)$, its degrees, and all non-zero
actions of matrix units on it.
\subpart Show that $\Bbbk e_{13}\subset Ae_{33}$ is a graded left submodule and
that $\bbD(M)\cong(Ae_{33}/\Bbbk e_{13})(1)$ as graded left $A$-modules.
\subpart Deduce that $\bbD(M)$ is not isomorphic to any
$(Ae_{jj})(r)$, $j\in\{1,2,3\}$, $r\in\mathbb Z$.
\end{subparts}

\emph{Hint for (iii).} Compare dimensions and the actions of $e_{11},e_{22},e_{33}$.
\end{parts}

\newpage
\titledquestion{Dependence on the symmetrising form}
Let $A=\kk A_2$, $S=\kk e_{11}\oplus\kk e_{22}$, $V=\kk e_{12}$, and
$M=e_{11}A$, with inherited gradings. For $c=(c_1,c_2)\in(\kk\setminus\{0\})^2$ define
\[
  \tau_c:S\to\kk,\qquad ae_{11}+be_{22}\mapsto c_1a+c_2b.
\]
Recall that $M^{\grstar}=\bigoplus_{d}\Hom_S(M_{-d},S)$ with the displayed
summand in degree $d$, that $\alpha_c:M^{\grstar}\to\bbD(M)$ is
$f\mapsto\tau_c\circ f$, and that
\[
  a\cdot_cf=\alpha_c^{-1}\bigl(a\,\alpha_c(f)\bigr),\qquad
  (a\lambda)(m)=\lambda(ma).
\]
A form $\tau$ on $S$ is \emph{symmetrising} if $\tau(st)=\tau(ts)$ for all
$s,t$ and the pairing $(s,t)\mapsto\tau(st)$ is non-degenerate.
\begin{parts}
\part Verify that $\tau_c$ is symmetrising. Let $\varepsilon:V\to\kk$ be
$\varepsilon(ue_{12})=u$. Find the right $S$-linear map $f:V\to S$ and the left
$S$-linear map $g:V\to S$ with $\tau_c\circ f=\varepsilon=\tau_c\circ g$, and
determine whether $f$ is also left linear and $g$ also right linear.

\part Let $\varphi_0(ae_{11}+be_{12})=ae_{11}$ and
$\varphi_1(ae_{11}+be_{12})=be_{22}$. Show that $\varphi_0,\varphi_1$ form a
homogeneous basis of $M^{\grstar}$ and find their degrees.

\part Compute the actions $e_{11}\cdot_c$, $e_{22}\cdot_c$, $e_{12}\cdot_c$ on
$\varphi_0,\varphi_1$ for general $c$. For $c,c'$, prove that
$\cdot_c=\cdot_{c'}$ if and only if $(c_1',c_2')$ is a scalar multiple of $(c_1,c_2)$.

\part \begin{subparts}
\subpart Explain, without computing any action, why
$(M^{\grstar},\cdot_c)\cong(M^{\grstar},\cdot_{c'})$ as graded left
$A$-modules for \emph{all} $c,c'$.
\subpart Write the isomorphism $t=\alpha_{c'}^{-1}\alpha_c$ explicitly on
$\varphi_0,\varphi_1$, and check directly that it commutes with
$e_{11},e_{22},e_{12}$.
\subpart Take $c=(1,1)$ and $c'=(1,2)$, $\operatorname{char}\kk\neq2$:
what is $t$?
\end{subparts}
\end{parts}

\end{questions}

\clearpage
\section*{Lecture 2: Tensor and quadratic algebras}
\begin{questions}
\setcounter{question}{2}

\titledquestion{Tensor products and $S$-duals}
Let $S$ be a finite-dimensional semisimple $\kk$-algebra. For a right
$S$-module $U$ let $U^*=\Hom_S(U,S)$; it is a left $S$-module by
$(sf)(u)=sf(u)$, and if $U$ is a bimodule then $(sft)(u)=sf(tu)$.
Let $V$ be a finite-dimensional right $S$-module and $W$ a finite-dimensional
$S$-bimodule, and define
\[
  \theta_r:W^*\otimes_SV^*\longrightarrow(V\otimes_SW)^*,\qquad
  \theta_r(g\otimes f)(v\otimes w)=g(f(v)w).
\]
\begin{parts}
\part Prove that $\theta_r$ is well defined: check bilinearity in $(g,f)$
and in $(v,w)$, invariance under $vs\otimes w\leftrightarrow v\otimes sw$ and
$gs\otimes f\leftrightarrow g\otimes sf$, and right $S$-linearity of
$\theta_r(g\otimes f)$. When $V$ is a bimodule, show
$\theta_r(sg\otimes ft)=s\,\theta_r(g\otimes f)\,t$.

\part Prove that $\theta_r$ is bijective.

\part Let $p:V\to V'$ be right $S$-linear and $q:W\to W'$ a bimodule map.
Prove $(p\otimes q)^*\circ\theta_r^{V',W'}=\theta_r^{V,W}\circ(q^*\otimes p^*)$.

\part Let $S=\kk^{\oplus3}$ with idempotents $e_1,e_2,e_3$, $V=\kk a$ with
$e_1ae_2=a$, and $W=\kk c$ with $e_2ce_3=c$. Compute
$W^*\otimes_SV^*$, $V^*\otimes_SW^*$ and $(V\otimes_SW)^*$, and describe
$\theta_r$. What happens if the two duals are tensored in the original order?

\part (Semisimplicity is needed.) Let $S=\kk[x]/(x^2)$ and $V=W=S/(x)$,
with the same definitions of $U^*$ and $\theta_r$. Show that
$W^*\otimes_SV^*$ and $(V\otimes_SW)^*$ are both $1$-dimensional but
$\theta_r=0$. At which step of your proof of (2) does the argument fail?
\end{parts}

\emph{Hint for (2).} First take $V=S$. Then use that every finite-dimensional
right $S$-module is a direct summand of $S^{\oplus n}$ for some $n$ (as a
right module; this is why $V$ is only assumed to be a right module).

\newpage
\titledquestion{Quadratic presentations and Hilbert series}
The parts are independent.
\begin{parts}
\part Let $n\geq3$, $Z\subset\{1,\dots,n-2\}$, and
$A_Z=\kk A_n/(e_{k,k+2}\mid k\in Z)$ with the quotient grading.
\begin{subparts}
\subpart Show that $\{e_{ii}\}\cup\{e_{ij}\mid i<j,\ \text{no }k\in Z\text{ with }i\leq k\leq j-2\}$
is a basis of $A_Z$.
\subpart Write $A_Z=T_S(V)/(R)$ with explicit $S$, $V$, and $R\subset V\otimes_SV$.
\subpart For $n=5$ and $Z=\{1,3\}$ compute $h_{A_Z}(t)$.
\subpart For $n=4$ and $Z=\{1,2\}$, show $e_{14}\in(e_{13},e_{24})$, that
$(A_Z)_2=0$, and that $A_Z=S\oplus V$ is radical-square-zero; compute $h_{A_Z}(t)$.
\end{subparts}
\emph{Hint.} $e_{ij}=e_{ik}e_{k,k+2}e_{k+2,j}$ whenever $i\leq k$ and $k+2\leq j$.

\part Let $S=\kk e_1\oplus\kk e_2$ and $V=\kk a\oplus\kk b$ with
$e_1ae_2=a$, $e_2be_1=b$ and all other idempotent actions zero, $V$ in degree~$1$.
Find a basis of $T_S(V)_d$ for all $d\geq0$ and compute $h_{T_S(V)}(t)$.
Then find a basis and $h$ for $T_S(V)/(ab)$. Which entries of the two
matrices differ, and why is $T_S(V)/(ab)$ finite-dimensional?

\emph{Hint.} Count alternating words in each degree, including degree $0$.

\part Let $V,W$ have bases $x_1,\dots,x_n$ and $y_1,\dots,y_m$, generators in
degree $1$, and
$\Sym(V)=T_\kk(V)/(x_ix_j-x_jx_i\mid i<j)$,
$\bigwedge W=T_\kk(W)/(w^2\mid w\in W)$.
\begin{subparts}
\subpart Compute $\dim R$ for the two relation spaces
(including linear independence) and deduce
$\dim\Sym(V)_2=\tfrac{n(n+1)}2$, $\dim(\bigwedge W)_2=\tfrac{m(m-1)}2$.
\subpart Using that ordered monomials form a basis of $\Sym(V)$ and strictly
increasing products a basis of $\bigwedge W$, compute
$\dim\Sym(V)_d$ and $\dim(\bigwedge W)_d$ for all $d$, and $h_{\Sym(V)}(t)$,
$h_{\bigwedge W}(t)$.
\subpart Let $P=\kk\langle x,y\rangle/(xy)$. Show that
$\{y^ax^b\mid a,b\geq0\}$ is a basis, and that $h_P(t)=h_{\kk[x,y]}(t)$.
Prove that $P$ and $\kk[x,y]$ are not isomorphic as algebras.
\end{subparts}
\end{parts}
\end{questions}

\clearpage
\section*{Lecture 3: The quadratic dual}

\begin{questions}
\setcounter{question}{4}

\titledquestion{Computing quadratic duals}
The parts are independent.
\begin{parts}
\part Let $A=\kk A_4/(e_{13})$.
Write a presentation of $A^!$ using reversed generators, and find homogeneous
bases and the Hilbert series of $A$ and $A^!$. Verify
\[
  h_A(t)\,h_{A^!}(-t)^{\mathsf T}=I_4,
\]
and check that dualising twice recovers $A$, identifying each twice-reversed
generator with its original matrix unit.

\emph{Hint.} Among $e_{12}e_{23}$ and $e_{23}e_{34}$ only the first is a
relation. Complement the composable pairs and reverse their order.

\part Let $S=\bigoplus_{i=1}^r\kk e_i$, let $V$ be a finite-dimensional
$S$-bimodule in degree $1$ with $C_{ij}=\dim_\kk e_iVe_j$, and let
$A=T_S(V)/(V\otimes_SV)=S\oplus V$.
\begin{subparts}
\subpart Prove $A^!=T_S(V^*)$, $h_A(t)=I+tC$, and
$h_{A^!}(t)=(I-tC^{\mathsf T})^{-1}$.
\subpart Deduce $h_A(t)h_{A^!}(-t)^{\mathsf T}=I$.
\subpart Specialise to $S=\kk e_1\oplus\kk e_2$, $V=\kk a\oplus\kk b$,
$e_1ae_2=a$, $e_2be_1=b$, and write down both series.
\end{subparts}

\emph{Hint.} For (i) use $V\otimes_SV$ as the relation space and the formula
$h_{T_S(V)}(t)=(I-tC)^{-1}$ of Lecture~2.

\part For $q\in\kk$ let
$C_q=\kk\langle x,y\rangle/(x^2,y^2,xy+qyx)$, $\deg x=\deg y=1$.
\begin{subparts}
\subpart Find a homogeneous basis of $C_q$ and compute all products of
two basis elements and $h_{C_q}(t)$, treating $q=0$ and $q\neq0$. Justify linear
independence of your basis.
\subpart For $u=\alpha x+\beta y$ compute $u^2$. For each $q$, determine
the $(\alpha,\beta)$ with $u^2=0$, and the $q$ for which every degree-$1$ element
squares to zero.
\subpart Prove: $C_q\cong C_{q'}$ as graded algebras if and only if either
$q=q'$ or $qq'=1$. Deduce that $C_1\not\cong C_{-1}$ when $\operatorname{char}\kk\neq2$
(they differ already as ungraded algebras, since one is commutative),
and say what changes in characteristic~$2$.
\subpart Compute a quadratic presentation of $C_q^!$ using $\theta_r$, and
its Hilbert series. Verify $h_{C_q}(t)h_{C_q^!}(-t)=1$. Compare $C_0^!$ with
$C_1^!$.
\end{subparts}

\emph{Hint for (3).} A graded isomorphism is determined by a linear bijection of
degree-$1$ parts and preserves the set of lines $\kk u$ with $u^2=0$.
For (4), evaluate a general combination of $X^2,XY,YX,Y^2$ on $x^2$, $y^2$,
and $xy+qyx$, with $X=x^*$, $Y=y^*$.
\end{parts}

\newpage
\titledquestion{Bases, Hilbert series, and quadratic duals}
The parts are independent.
\begin{parts}
\part \textbf{A square with a parameter.}
Let $S=\kk e_1\oplus\kk e_2\oplus\kk e_3\oplus\kk e_4$ and let $V$ have basis
$a,b,c,d$ with
\[
  e_1ae_2=a,\quad e_2be_4=b,\quad e_1ce_3=c,\quad e_3de_4=d,
\]
and all other idempotent actions zero. The only composable words of length $2$
are $ab$ and $cd$. For $\lambda\in\kk$ set $A_\lambda=T_S(V)/(ab-\lambda cd)$.
\begin{subparts}
\subpart For every $\lambda$, including $0$, find a presentation of
$A_\lambda^!$, bases of $A_\lambda$ and $A_\lambda^!$, and both Hilbert
series. Do they depend on $\lambda$? Dualise again to recover the relation.
\subpart Prove that $A_\lambda\cong A_\mu$ as graded algebras over $S$
if and only if $\lambda=\mu=0$ or both are non-zero.
\subpart Show $A_1^!\cong(A_1)^{\op}$ as graded algebras over $S$.
\end{subparts}
\emph{Hint for (1).} Evaluate $\mu b^*a^*+\nu d^*c^*$ on $ab-\lambda cd$.
For (2), a graded isomorphism over $S$ sends each of $a,b,c,d$ to a scalar multiple
of itself (or swaps the two paths); for (3) rescale $d^*$.

\part \textbf{Counting words.} Let $B=T_\kk(\kk x\oplus\kk y)/(x^2)$,
$\deg x=\deg y=1$.
\begin{subparts}
\subpart Describe a basis of $B_d$, and compute $a_d=\dim B_d$ for $0\leq d\leq6$.
\subpart Express $a_d$ through $a_{d-1},a_{d-2}$ by sorting words by first
letter, and find $h_B(t)$ as a rational function.
\subpart Compute $B^!$ (with $X=x^*$, $Y=y^*$), a homogeneous basis, and
$h_{B^!}(t)$. Verify $h_B(t)h_{B^!}(-t)=1$.
\end{subparts}
\end{parts}
\end{questions}

\end{document}
